% concatenated from Top_Star.tex
\documentclass[10pt,a4paper]{article}
\usepackage{fullpage}
\usepackage{amsfonts,amsmath,amssymb}
\usepackage{amsthm}
\usepackage{graphicx}\usepackage{sectsty}
\usepackage{authblk}
%------------------------------------------------------------------------%

\theoremstyle{plain}
\newtheorem{theorem}{Theorem}[section]
\newtheorem{lemma}[theorem]{Lemma}
\newtheorem{proposition}[theorem]{Proposition}
\newtheorem{corollary}[theorem]{Corollary}

\theoremstyle{definition}
\newtheorem{definition}[theorem]{Definition}
\theoremstyle{remark}
\newtheorem{remark}[theorem]{Remark}
%\newenvironment{proof}[1][Proof]{\begin{trivlist}
%\item[\hskip \labelsep {\bfseries #1}]}{\end{trivlist}}
%\newenvironment{definition}[1][Definition]{\begin{trivlist}
%\item[\hskip \labelsep {\bfseries #1}]}{\end{trivlist}}
\newenvironment{example}[1][Example]{\begin{trivlist}
\item[\hskip \labelsep {\bfseries #1}]}{\end{trivlitst}}


%\newtheorem{remark}[theorem]{Remark}
%\newenvironment{remark}[1][Remark]{\begin{trivlist}
%\item[\hskip \labelsep {\bfseries #1}]}{\end{trivlist}}
%%%???? ????

\numberwithin{equation}{section}
\def\address#1{\begin{center}{\it #1}\end{center}}
\def\email#1{\begin{center}{\it #1}\end{center}}
\sectionfont{\large}
\newenvironment{acknowledgement}[1][Acknowledgement
]{\begin{trivlist} \item[\hskip \labelsep {\bfseries
#1}]}{\end{trivlist}}

\begin{document}
\title{Recovering the Topology in One Point Interaction Problem on Extended Non-Local Star Graphs}
\author[1]{Lung-Hui Chen}
\affil[1]{\footnotesize General Education Center, Ming Chi University of Technology, New Taipei City, 24301, Taiwan.}
\maketitle
\begin{abstract}
The author studies the inverse spectral problem of Sturm-Liouville operator on a star-like graph. To this star-like graph centered at the origin as its vertex, there are attached $m$ edges that imposed the Sturm-Liouville operator with certain non-local potential functions with some suitable local boundary value conditions. At the vertex, we consider one point interaction condition at vertex to model a network that fixed on the end of the edges on the graph. The vibration and flow changes are monitored at that vertex which serves as certain control/regulation center. The author shows that the system is solvable under very necessary conditions. It is crucial to recover  the topology of the network. In this paper, author constructs the special solution edge by edge and point to point.
\\MSC: 47A55/34A55/34K29.
\\Keywords:  Fourier analysis; star graph; inverse spectral problem; network theory; non-local regulator; boundary control problem; traffic control.
\end{abstract}
\section{Introduction}
This work explores the inverse spectral analysis of Sturm-Liouville operators situated on star-shaped metric graphs. On such graphs, where multiple edges emanate from a central vertex, provide an effective mathematical model for analyzing signal transmission and dynamic processes within structured networks. The application of Sturm-Liouville theory in this context has become increasingly relevant in bridging mathematical physics with real-world engineering systems such as telecommunications, energy grids, and transport logistics, as noted in  studies \cite{Freiling,Kostrykin,Kuchment,Kurasov,Novikov,Pokornyi}.

The behavior of wave propagation and diffusion on these networks can be described by differential operators acting along each edge, governed by local boundary conditions and damped by non-local potential terms. At the central vertex, a fixed-point (frozen argument) condition is imposed, simulating a centralized monitoring or control mechanism. This setting allows us to model systems where information or energy must flow efficiently and reversibly across the distributed structure.

The non-local boundary conditions are not merely technical features, but also contribute to the stability and physical realization of certain engineering modeling. In particular, they help explain the physical principles such as the conservation of energy which in turn the mathematical uniqueness of solutions of the state. The investigation presented here demonstrates that, under appropriate local or non-local conditions, the inverse spectral problem admits an inverse unique and solvable formulation.


\par
The most interesting aspect of such modeling is its connection to the theory of quantum graphs, which studies the physical dynamics on metric graphs governed by various sorts of differential operators \cite{Alb,Bon,Freiling,Gerasimenko,Kottos,Naimark}, along with many different kinds of boundary conditions which often compromise the energy conservation laws or even certain commercial requirements in real world scenarios. The boundary condition acts as an integral part of the modeling these specific problems. It plays a role in many engineering projects. We are interested to formulate and to solve a few number of direct and inverse spectral problems and scattering problems on quantum graphs \cite{Belishev,Carlson,Kostrykin,Kurasov,Kurasov2,Pivovarchik1,Yurko1,Yurko2}.
Specifically, in this paper, we study the inverse spectral problems for a Schr\"{o}dinger operator with non-local potential functions on a star graph. Previously, the inverse spectral problems on a finite graph for the Sturm-Liouville problems with non-local potentials or with various boundary value conditions were considered in \cite{Alb,Alb2,Bon,Nizhnik1,Nizhnik2,Pivovarchik1,Yang}.
\par
Let us consider the following star graph $T$: The central vertex of $T$ is located at the origin $0\in\mathbb{R}^{2}$ and connected by 
$m$ edges, each sharing the origin as their common vertex, which we denoted as $V_{0}$.  Each edge $e_{j}$ has length $l_{j}$, $j=1,2,\ldots,m$ with end points denoted as $v_{j}$, $j=1,2,\ldots,m$.
    \par
On each edge $e_{j}$, we assume that there exists the function $\psi_{j}(x)\in W_{2}^{2}(0,l_{j})$ such that the function $\psi_{j}(x)$ satisfies the following eigenvalue problem with complex-valued non-local potentials $q_{j}(x)\in L^{2}(0,l_{j})$:
\begin{eqnarray}\label{1.1}
-\frac{d^{2}}{dx^{2}}\psi_{j}(x)+q_{j}(x)\psi_{j}(0)=\lambda\psi_{j}(x),\,0<x<l_{j},\,j=1,2,\ldots,m,
\end{eqnarray}
in which $m\geq2$, $\lambda\in\mathbb{C}$ is the spectral parameter,
and it is imposed with the boundary conditions
\begin{equation}\label{1.2}
\psi_{j}(l_{j})=0,\,j=1,2,\ldots,m;\,\psi_{1}(0)=\psi_{2}(0)=\cdots=\psi_{m}(0),
\end{equation}
and the non-local boundary value condition
\begin{equation}\label{1.3}
\sum_{j=1}^{m}\big\{\frac{d}{dx}\psi_{j}(0)-\int_{0}^{l_{j}}\psi_{j}(x)\overline{q_{j}(x)}dx\big\}=0.
\end{equation}
Also, $\lambda=(\frac{z\pi}{l_{j}})^{2}$, $z\in\mathbb{C}$, and the solution $\psi_{j}(x,z)$ depends on $(x,z)\in\mathbb{R}\times\mathbb{C}$ for each $j$.
We note that the inverse spectral problems is posed with a frozen argument at the vertex that is called a non-local point interaction in quantum mechanics, and  the equation actually reads as
\begin{equation}\label{1.4}
\sum_{j=1}^{m}\big\{\frac{d}{dx}\psi_{j}(0,z)-\int_{0}^{l_{j}}\psi_{j}(x,z)\overline{q_{j}(x)}dx\big\}=0,\,z\in\mathbb{C}.
\end{equation}
In this paper, we use the spectral data to recover the information on angular information between the edges $$\{e_{1},e_{2},\cdots,e_{m}\}.$$
To this end, we consider the extended closed graph from the star graph $T$, which is denoted as $\overline{T}$ that is obtained by connecting the end point $v_{j}$, $j=1,2,\ldots,m$, and $v_{m+1}=v_{1}$. Let us denote the angle between $e_{j}$ and $e_{j+1}$ as $\theta_{j}$, and the angle between $e_{m}$ and $e_{1}$ as $\theta_{m}$. Moreover, we call the new edges between the vertices $v_{j}$ and $v_{j+1}$ as $\bar{e}_{j}$ with length $\bar{l}_{j}$, and the edge between $v_{m}$ and $v_{1}$ as $\bar{e}_{m}$ with length $\bar{l}_{m}$. 
We set $\cup_{j=1}^{m}\bar{e}_{j}$ to be a cyclic metric graph by taking $v_{m+1}=v_{1}$.
Surely, $v_{j},\,j\geq1,$ is a vertex with edge $e_{j}$, $\bar{e}_{j-1}$, and $\bar{e}_{j+1}$. The extension graph $\overline{T}$ is trivially unique.
\par
Therefore, we modify the differential system~(\ref{1.1}),~(\ref{1.2}), and~(\ref{1.4}) on the presence of  $\bar{e}_{j'}$, $j'=1,2,\cdots,m.$ For $z\in\mathbb{C}$,
\begin{eqnarray}\label{1.5}
\left\{
  \begin{array}{llllll}\vspace{7pt}
-\frac{d^{2}}{dx^{2}}\psi_{j}(x,z)+q_{j}(x)\psi_{j}(0,z)=(\frac{z\pi}{l_{j}})^{2}\psi_{j}(x,z),\,x\in e_{j},\,j=1,2,\ldots,m;\\\vspace{7pt}
-\frac{d^{2}}{dx^{2}}\psi_{j'}(x,z)=(\frac{z\pi}{l_{j}})^{2}\psi_{j'}(x,z),\,x\in \bar{e}_{j'},\,j'=1,2,\ldots,m;\\\vspace{7pt}
\psi_{j}(l_{j},z)=0,\,j=1,2,\ldots,m;\\\vspace{7pt}
\psi_{1}(0,z)=\psi_{2}(0,z)=\cdots=\psi_{m}(0,z);\\\vspace{7pt}
\psi_{j'}(l_{j'}^{-},z)=\psi_{j'-1}(\bar{l}_{j'-1}^{-},z)=\psi_{j'}(0^{+},z)=0,\,j=1,\cdots,m;\\\vspace{7pt}
\sum_{j=1}^{m}\sum_{j'=1}^{m}\big\{\frac{d}{dx}\psi_{j}(0^{+},z)+\frac{d}{dx}\psi_{j'}(l_{j'}^{-},z)+\frac{d}{dx}\psi_{j'-1}(\bar{l}_{j'-1}^{-},z)+\frac{d}{dx}\psi_{j'}(0^{+},z)\\\hspace{80pt}-\int_{0}^{l_{j}}\psi_{j}(x,z)\overline{q_{j}(x)}dx\big\}=0,
  \end{array}
\right.
\end{eqnarray}
in which $\frac{d}{dx}\psi_{j'}(l_{j'}^{-},z)+\frac{d}{dx}\psi_{j'-1}(\bar{l}_{j'-1}^{-},z)+\frac{d}{dx}\psi_{j'}(0^{+},z)$ prepares for some kind Kirchhoff's condition at vertex $v_{j'}$ on $\cup_{j=1}^{m}\bar{e}_{j}$ and the last line is the non-local Kirchhoff's condition in the system which collects all the eigenvalues of~(\ref{1.5}). When using subscript $j'$, we mean to count the objects on $\cup_{j=1}^{m}\bar{e}_{j}$.
In this paper, we aim to recover the information on $\theta_{j}$, $j=1,2,\ldots,m$, while the knowledge of $l_{j}$ and $q_{j}(x)$, $j=1,2,\ldots,m$, is given.



\section{Special Solution and Characteristic Function}
In this section, we construct the special solutions for system~(\ref{1.5}).
We start with the Fourier-sine series in $L^{2}(0,l_{j})$:
$$q_{j}(x)=\sum_{n=1}^{\infty}q_{j,n}\sin[\frac{n\pi}{l_{j}} (l_{j}-x)]$$ with Fourier-Sine transform and its coefficients
\begin{eqnarray}\label{2.1}
&&\mathcal{F}(q_{j})(z)=\frac{2}{l_{j}}\int_{0}^{l_{j}}q_{j}(x)\sin[\frac{z\pi}{l_{j}} (l_{j}-x)]dx;\\
&&q_{j,n}=\frac{2}{l_{j}}\int_{0}^{l_{j}}q_{j}(x)\sin[\frac{n\pi}{l_{j}} (l_{j}-x)]dx.
\end{eqnarray}
Firstly, we consider a special solution of system~(\ref{1.1}) with $\lambda=(\frac{z\pi}{l_{j}})^{2}$ to meet all conditions in~(\ref{1.1}). We choose
\begin{equation} \label{2.2}
\phi_{j}(x;z)=\Big(\sin{z\pi(1-\frac{x}{l_{j}})}+\sin{z l_{j}}\sum_{n=1}^{\infty}\frac{q_{j,n}\sin[\frac{n\pi}{l_{j}}(l_{j}-x)]}{z^{2}-(\frac{n\pi}{l_{j}})^{2}}\Big)\prod_{k\neq j}\sin{z l_{k}},\,x\in T,
\end{equation} and then we can verify
\begin{equation}
\phi_{1}(0;z)=\phi_{2}(0;z)=\cdots=\phi_{m}(0;z)=\sin z\pi;
\end{equation}
\begin{equation}
\phi_{1}(l_{1};z)=\phi_{2}(l_{2};z)=\cdots=\phi_{m}(l_{m};z)=0.
\end{equation}
Thus, one point interaction at $v_{0}=0$ in~(\ref{1.5}) is satisfied.
Moreover,
\begin{equation}\label{3333}
\phi_{j}'(x;z)=\Big(-\frac{z\pi}{l_{j}}\cos{z\pi(1-\frac{x}{l_{j}})}-\sin{z l_{j}}\sum_{n=1}^{\infty}\frac{n\pi}{l_{j}}\frac{q_{j,n}\cos[\frac{n\pi}{l_{j}}(l_{j}-x)]}{z^{2}-(\frac{n\pi}{l_{j}})^{2}}\Big)\prod_{k\neq j}\sin{z l_{k}},\,x\in T,
\end{equation} 
and then
\begin{eqnarray}\label{2.3}
\phi_{j}'(0;z)=\Big(-\frac{z\pi}{l_{j}}\cos{z\pi}-\sin{z l_{j}}\sum_{n=1}^{\infty}\frac{(-1)^{n}n\pi}{l_{j}}\frac{q_{j,n}}{z^{2}-(\frac{n\pi}{l_{j}})^{2}}\Big)\prod_{k\neq j}\sin{z l_{k}},\,j=1,2,\ldots,m.
\end{eqnarray}
To prepare a Kirchhoff condition at $x=v_{0}$, we sum over $j$-th summands.
\begin{eqnarray}
\sum_{j=1}^{m}\phi_{j}'(0;z)=-\sum_{j=1}^{m}\Big(\frac{z\pi}{l_{j}}\cos{z\pi}+\sin{z l_{j}}\sum_{n=1}^{\infty}\frac{(-1)^{n}n\pi}{l_{j}}\frac{q_{j,n}}{z^{2}-(\frac{n\pi}{l_{j}})^{2}}\Big)\prod_{k\neq j}\sin{z l_{k}}.\label{25}
\end{eqnarray}
\par
Due to the presence of $\{\bar{e}_{j}\}_{j=1}^{m}$ and $\{v_{j}\}_{j=1}^{m}$ on $\overline{T}$, we add up the special solutions on each $\bar{e}_{j}$:
\begin{equation}
\phi_{j}(x;z)=\sin{z\pi}(1-\frac{x}{\bar{l}_{j}})\prod_{k\neq j}\sin{z l_{k}},   0<x\leq\bar{l}_{j},\label{555}
\end{equation}
in which we consider the parametrization of $x$ as its are length when $x$ moves along $\cup_{j=1}^{m}\bar{e}_{j}$ that is a closed cyclic graph, and $v_{m+1}=v_{1}$. Therefore, we define $\phi_{j}(0;z):=\phi_{j-1}(\bar{l}_{j-1};z)=0$ which is zero due to~(\ref{555}).
The zero set of $\sin{z\pi}(1-\frac{x}{\bar{l}_{j}})$ is $x\in\{\frac{\bar{l}_{j}(z-n)}{z}\}_{j=1}^{m}$ and $z\in\mathbb{Z}$. Most importantly, we obtain the continuity condition at $\{v_{j}\}_{j=1}^{m}$.
Hence,
\begin{equation}
\phi_{j}(0;z)=\phi_{j+1}(0;z)=0.
\end{equation}
We also observe $$-\phi_{j}''(x;z)=(\frac{z\pi}{\bar{l}_{j}})^{2}\phi_{j}(x;z),\,x\in \bar{e}_{j}.$$ The Dirichlet problem has the general solution. To prepare the Kirchhoff condition for $x$ at vertex $v^{j}$, $j=1,2,\cdots,m,$
we use~(\ref{3333}) and~(\ref{555}) to deduce that
\begin{eqnarray}\nonumber
&&\hspace{12pt}\frac{d}{dx}\phi_{j'}(l_{j'}^{-},z)+\frac{d}{dx}\phi_{j'-1}(\bar{l}_{j'-1}^{-},z)+\frac{d}{dx}\phi_{j'}(0^{+},z)\\
&&\nonumber=-\Big(\frac{z\pi}{l_{j'}}+\frac{z\pi}{\bar{l}_{j'-1}}+\frac{z\pi}{\bar{l}_{j'}}\cos{z\pi}+\sin{zl_{j}}\sum_{n=1}^{\infty}\frac{n\pi}{l_{j}}\frac{q_{j,n}}{z^{2}-(\frac{n\pi}{l_{j}})^{2}}\Big)\prod_{k\neq j}\sin{z l_{k}},\,j=1,2,\ldots,m.\label{288}%%%%%%%%%%unfinished
\end{eqnarray}
Secondly, we recall the non-local condition in~(\ref{1.5}) on star graph $T$,
\begin{eqnarray}\nonumber
&&\int_{0}^{l_{j}}\phi_{j}(x;z)\overline{q_{j}(x)}dx\\&=&\Big(\int_{0}^{l_{j}}\sin{z\pi(1-\frac{x}{l_{j}})}\overline{q_{j}(x)}dx+\sin{z l_{j}}\sum_{n=1}^{\infty}\frac{q_{j,n}\int_{0}^{l_{j}}\sin[\frac{n\pi}{l_{j}}(l_{j}-x)]\overline{q_{j}(x)}dx}{z^{2}-(\frac{n\pi}{l_{j}})^{2}}\Big)\prod_{k\neq j}\sin{z l_{k}}.
\end{eqnarray}
Moreover, on the edges $\{e_{j}\}$, we consider the summation
\begin{eqnarray}\nonumber
&&\sum_{j=1}^{m}\int_{0}^{l_{j}}\phi_{j}(x;z)\overline{q_{j}(x)}dx\\&=&\sum_{j=1}^{m}\Big(\int_{0}^{l_{j}}\sin{z\pi(1-\frac{x}{l_{j}})}\overline{q_{j}(x)}dx+\sin{z l_{j}}\sum_{n=1}^{\infty}\frac{q_{j,n}\int_{0}^{l_{j}}\sin[\frac{n\pi}{l_{j}}(l_{j}-x)]\overline{q_{j}(x)}dx}{z^{2}-(\frac{n\pi}{l_{j}})^{2}}\Big)\prod_{n\neq j}\sin{z l_{k}}.\label{29}
\end{eqnarray}
On the extended edges $\{\bar{e}_{j}\}$, $\overline{q_{j}(x)}\equiv 0$. Therefore, the equation~(\ref{2.3}),~(\ref{25}),~(\ref{288}), and~(\ref{29}) add up to be the characteristic function in this paper:
\begin{eqnarray}\nonumber
\Phi(z)&\hspace{-5pt}:=\hspace{-5pt}&\sum_{j=1}^{m}\Big(\int_{0}^{l_{j}}\sin{z\pi(1-\frac{x}{l_{j}})}\overline{q_{j}(x)}dx+\sin{z l_{j}}\sum_{n=1}^{\infty}\frac{q_{j,n}\overline{q_{j,n}}}{z^{2}-(\frac{n\pi}{l_{j}})^{2}}\Big)\prod_{k\neq j}\sin{z l_{k}}\\&&
+\sum_{j=1}^{m}\Big(\frac{z\pi}{l_{j}}\cos{z\pi}+\sin{z l_{j}}\sum_{n=1}^{\infty}\frac{(-1)^{n}n\pi}{l_{j}}\frac{q_{j,n}}{z^{2}-(\frac{n\pi}{l_{j}})^{2}}\Big)\prod_{k\neq j}\sin{z l_{k}}\nonumber\\
&&+\sum_{j'=1}^{m}\Big(\frac{z\pi}{l_{j'}}+\frac{z\pi}{\bar{l}_{j'-1}}+\frac{z\pi}{\bar{l}_{j'}}\cos{z\pi}+\sin{zl_{j'}}\sum_{n=1}^{\infty}\frac{n\pi}{l_{j'}}\frac{q_{j,n}}{z^{2}-(\frac{n\pi}{l_{j'}})^{2}}\Big)\prod_{k\neq j'}\sin{z l_{k}},\label{Phi}
\end{eqnarray}
in which we set $\bar{l}_{0}=\bar{l}_{m}$.


\section{Results}
We need to complete the system~(\ref{1.5}) firstly.
\begin{lemma}
The set $\mathbb{Z}$ satisfied the first condition of~(\ref{1.5}).
\end{lemma}
\begin{proof}
It suffices to verify the first condition as~(\ref{1.1}). We proceed with~(\ref{2.2}) to obtain
\begin{equation}
-\phi''_{j}(x;z)=\Big((\frac{z\pi}{l_{j}})^{2}\sin{z\pi(1-\frac{x}{l_{j}})}+(\frac{n\pi}{l_{j}})^{2}\sin{z l_{j}}\sum_{n=1}^{\infty}\frac{q_{j,n}\sin[\frac{n\pi}{l_{j}}(l_{j}-x)]}{z^{2}-(\frac{n\pi}{l_{j}})^{2}}\Big)\prod_{k\neq j}\sin{z l_{k}},\label{331}
\end{equation}
\begin{equation}
q_{j}(x)\phi_{j}(0;z)=q_{j}(x)\Big(\sin{z\pi}+\sin{z l_{j}}\sum_{n=1}^{\infty}\frac{q_{j,n}\sin[n\pi]}{z^{2}-(\frac{n\pi}{l_{j}})^{2}}\Big)\prod_{k\neq j}\sin{z l_{k}},\label{332}
\end{equation}
and
\begin{equation}
(\frac{z\pi}{l_{j}})^{2}\phi_{j}(x;z)=(\frac{z\pi}{l_{j}})^{2}\Big(\sin{z\pi(1-\frac{x}{l_{j}})}+\sin{z l_{j}}\sum_{n=1}^{\infty}\frac{q_{j,n}\sin[\frac{n\pi}{l_{j}}(l_{j}-x)]}{z^{2}-(\frac{n\pi}{l_{j}})^{2}}\Big)\prod_{k\neq j}\sin{z l_{k}}.\label{333}
\end{equation}
Most importantly, for $z\in\mathbb{Z}$, we obtain
$$-\phi''_{j}(x;z)+q_{j}(x)\phi_{j}(0;z)=(\frac{z\pi}{l_{j}})^{2}\phi_{j}(x;z).$$
Then,~(\ref{1.1}) is satisfied on $T$. 
\end{proof}%%%%%Not finished!!How about on cyclic graph? only for integers. Becareful!!



\begin{proposition}\label{T3.2}
Let the differential system~(\ref{1.5}) be defined on the star graph 
$\overline{T}^{\sigma}$, $\sigma=1,2,$ with the identical potential function on each  edge $e_{j}$ with length $l_{j}$, $1\leq j\leq m$. Let $\Phi^{\sigma}(z)$ be the characteristic defined on $\overline{T}^{\sigma}$, $\sigma=1,2,$ If $ \Phi^{1}(z)=\Phi^{2}(z)$ for $z\in\mathbb{Z}$ for all $j$'s, then $\theta^{1}_{j}=\theta^{2}_{j}$ for $1\leq j\leq m$.
\end{proposition}
\begin{proof}
Under the assumption of proposition, we derive from~(\ref{Phi}) that
\begin{equation}
\Phi^{1}(z)-\Phi^{2}(z)=\sum_{j'=1}^{m}
\{\frac{z\pi}{\bar{l}_{j'-1}^{1}}-\frac{z\pi}{\bar{l}_{j'-1}^{2}}+\frac{z\pi}{\bar{l}_{j'}^{1}}\cos{z\pi}-\frac{z\pi}{\bar{l}_{j'}^{2}}\cos{z\pi}\}\prod_{k\neq j'}\sin{z l_{k}}=0,\,z\in\mathbb{Z}.\label{3.4}
\end{equation}%%%%%%%%%%need rationally independence????
Most importantly, we set
\begin{equation}
\mathcal{Z}_{j'}:=\big\{z|\,\sin{z\pi}\prod_{k\neq j'}\sin{z l_{k}}=0\big\}.
\end{equation}
For each $j'=1,2,\cdots,m$, we plug in $\mathcal{Z}_{j'}$ consecutively into~(\ref{3.4}) to obtain
\begin{equation}
\frac{1}{\bar{l}_{j'-1}^{1}}-\frac{1}{\bar{l}_{j'-1}^{2}}+\frac{1}{\bar{l}_{j'}^{1}}(-1)^{n}-\frac{1}{\bar{l}_{j'}^{2}}(-1)^{n}=0,\,z\in\mathcal{Z}_{j'}.\label{3.6}
\end{equation}
Because $n$ can be either odd or even integers, we have shown $\bar{l}^{1}_{j'-1}=\bar{l}^{2}_{j'-1}$, $1\leq j'\leq m$.
Thus,
we prove that that all extended edges are equal from both sides. This is the uniqueness of $\{\theta_{j}\}_{j=1}^{m}$.

\end{proof}
One step further, if the knowledge of $\{l_{j}\}_{j=1}^{m}$ and $\{\bar{l}_{j}\}_{j=1}^{m}$ are given, may we recover the information on $\{q_{j}\}_{j=1}^{m}$?
\begin{lemma}\label{LL33}
The following identities hold.
\begin{equation}
\sin{z l_{j}}\sum_{n=1}^{\infty}\frac{(-1)^{n}n\pi}{l_{j}}\frac{q_{j,n}}{z^{2}-(\frac{n\pi}{l_{j}})^{2}}=\int_{0}^{\pi}q(\frac{l_{j}}{\pi}(\pi-x)\sin{zx}dx.
\end{equation}
If $q_{j,n}=0$, $n\in\mathbb{Z}$, then $q_{j}(x)\equiv0$, $0\leq j\leq m$.
\end{lemma}
%%%%%%%%%%%%%%%%%%%%%%%%%%%match with (2.2)???????
\begin{proof}
We recall that 
$$q_{j}(x)=\sum_{n=1}^{\infty}q_{j,n}\sin[\frac{n\pi}{l_{j}} (l_{j}-x)]$$ with Fourier coefficients
$$q_{j,n}=\frac{2}{l_{j}}\int_{0}^{l_{j}}q_{j}(x)\sin[\frac{n\pi}{l_{j}} (l_{j}-x)]dx.$$
Using
\begin{eqnarray}
\frac{q_{j,n}}{z^{2}-(\frac{n\pi}{l_{j}})^{2}}=\frac{-\frac{1}{2}\frac{l_{j}}{n\pi}q_{j,n}}{(z+\frac{n\pi}{l_{j}})}+\frac{\frac{1}{2}\frac{l_{j}}{n\pi}q_{j,n}}{(z-\frac{n\pi}{l_{j}})},
\end{eqnarray}
we deduce
\begin{eqnarray}\nonumber
\sin{z l_{j}}\sum_{n=1}^{\infty}\frac{(-1)^{n}n\pi}{l_{j}}\frac{q_{j,n}}{z^{2}-(\frac{n\pi}{l_{j}})^{2}}&\hspace{-4pt}=\hspace{-4pt}&\sin{z l_{j}}\sum_{n=1}^{\infty}(-1)^{n}\frac{-\frac{1}{2}q_{j,n}}{(z+\frac{n\pi}{l_{j}})}+\sin{z l_{j}}\sum_{n=1}^{\infty}(-1)^{n}\frac{\frac{1}{2}q_{j,n}}{(z-\frac{n\pi}{l_{j}})}\\\nonumber
&\hspace{-4pt}=\hspace{-4pt}&\sin{z l_{j}}\sum_{n=-1}^{\infty}(-1)^{n}{l_{j}}\frac{\frac{1}{2}q_{j,n}}{(z-\frac{n\pi}{l_{j}})}+\sin{z l_{j}}\sum_{n=1}^{\infty}(-1)^{n}\frac{\frac{1}{2}q_{j,n}}{(z-\frac{n\pi}{l_{j}})}
\\&\hspace{-4pt}=\hspace{-4pt}&\sin{z l_{j}}\sum_{n=-\infty}^{\infty}(-1)^{n}\frac{\frac{1}{2}q_{j,n}}{(z-\frac{n\pi}{l_{j}})}\nonumber
\\&\hspace{-4pt}=\hspace{-4pt}&\sin{\eta\pi}\sum_{n=-\infty}^{\infty}\frac{(-1)^{n}}{\pi}\frac{\frac{1}{2\pi}q_{j,n}}{(\frac{zl_{j}}{\pi}-n)}.
\end{eqnarray}
Let $z=\frac{\eta\pi}{l_{j}}$, $\eta\in\mathbb{C}$, and $x=\frac{l_{j}}{\pi}y$. Therefore,
\begin{eqnarray}\nonumber
&&\sin{z l_{j}}\sum_{n=1}^{\infty}\frac{(-1)^{n}n\pi}{l_{j}}\frac{q_{j,n}}{z^{2}-(\frac{n\pi}{l_{j}})^{2}}\\&\hspace{-4pt}=&\sin{\eta\pi}\sum_{n=-\infty}^{\infty}\frac{(-1)^{n}}{\pi}\frac{\frac{1}{2\pi}q_{j,n}}{(\frac{zl_{j}}{\pi}-n)}\nonumber
\\&\hspace{-4pt}=&\sin{\eta\pi}\sum_{n=-\infty}^{\infty}(-1)^{n}\frac{\frac{1}{2\pi}\int_{0}^{\pi}q(\frac{l_{j}}{\pi}(\pi-x)\sin{nx}dx}{\pi(\eta-n)},\label{3100}
\end{eqnarray}
which is a Fourier sine transform as we apply the interpolation theory in complex analysis referring to M. Cartwright and Levin \cite[p.\,150,\,151]{Levin2}, and
$\{\frac{1}{2\pi}\int_{0}^{\pi}q(\frac{l_{j}}{\pi}(\pi-x)\sin{nx}dx\}_{n}\in l^{2}.$
A Fourier transform vanishes at $\mathbb{Z}$ must be trivial \cite[p.\,150,\,Theorem\,1]{Levin2}. The lemma is proven.

\end{proof}
\begin{theorem}\label{T3.3}
Let the differential system~(\ref{1.5}) be defined on the star graph 
$\overline{T}^{\sigma}$, $\sigma=1,2,$ with the potential function $q^{\sigma}_{j}$ defined on each  edge $e_{j}$ with length $l_{j}$, $1\leq j\leq m$. Let $\Phi^{\sigma}(z)$ be the characteristic defined on $\overline{T}^{\sigma}$, $\sigma=1,2.$ If $ \Phi^{1}(z)=\Phi^{2}(z)$ for $z\in\frac{\mathbb{Z}\pi}{l_{j}}$, then $q^{1}_{j}=q^{2}_{j}$ for $1\leq j\leq m$.
\end{theorem}
\begin{proof}
We consider the subtraction from~(\ref{Phi}).
\begin{eqnarray}\nonumber
\Phi^{1}(z)-\Phi^{2}(z)&\hspace{-5pt}=\hspace{-5pt}&\sum_{j=1}^{m}\Big(\int_{0}^{l_{j}}\sin{z\pi(1-\frac{x}{l_{j}})}\overline{q^{1}_{j}(x)-q^{2}_{j}(x)}dx+\sin{z l_{j}}\sum_{n=1}^{\infty}\frac{q^{1}_{j,n}\overline{q^{1}_{j,n}}-q^{2}_{j,n}\overline{q^{2}_{j,n}}}{z^{2}-(\frac{n\pi}{l_{j}})^{2}}\Big)\prod_{k\neq j}\sin{z l_{k}}\\&&
+\sum_{j=1}^{m}\Big(\sin{z l_{j}}\sum_{n=1}^{\infty}\frac{(-1)^{n}n\pi}{l_{j}}\frac{q^{1}_{j,n}-q^{2}_{j,n}}{z^{2}-(\frac{n\pi}{l_{j}})^{2}}\Big)\prod_{k\neq j}\sin{z l_{k}}\nonumber\\
&&+\sum_{j=1}^{m}\Big(\sin{zl_{j}}\sum_{n=1}^{\infty}\frac{n\pi}{l_{j}}\frac{q^{1}_{j,n}-q^{2}_{j,n}}{z^{2}-(\frac{n\pi}{l_{j}})^{2}}\Big)\prod_{k\neq j}\sin{z l_{k}}.\label{3.8}
\end{eqnarray}
Now we set%%%%%%%%%%%need rational independence?
\begin{equation}
\mathbf{Z}_{j}:=\{z|\,\sin{z l_{j}}\prod_{k\neq j}\sin{z l_{k}}=0\}.
\end{equation}
For each $j$, let us plug in $\mathbf{Z}_{j}$ into~(\ref{3.8}) consecutively under the assumption of theorem  to deduce
\begin{eqnarray}\nonumber
\Phi^{1}(z)-\Phi^{2}(z)&\hspace{-5pt}=\hspace{-5pt}&\int_{0}^{l_{j}}\sin{z\pi(1-\frac{x}{l_{j}})}\overline{q^{1}_{j}(x)-q^{2}_{j}(x)}dx+\sin{z l_{j}}\sum_{n=1}^{\infty}\frac{q^{1}_{j,n}\overline{q^{1}_{j,n}}-q^{2}_{j,n}\overline{q^{2}_{j,n}}}{z^{2}-(\frac{n\pi}{l_{j}})^{2}}\\&&
+\sin{z l_{j}}\sum_{n=1}^{\infty}\frac{(-1)^{n}n\pi}{l_{j}}\frac{q^{1}_{j,n}-q^{2}_{j,n}}{z^{2}-(\frac{n\pi}{l_{j}})^{2}}+\sin{zl_{j}}\sum_{n=1}^{\infty}\frac{n\pi}{l_{j}}\frac{q^{1}_{j,n}-q^{2}_{j,n}}{z^{2}-(\frac{n\pi}{l_{j}})^{2}}=0,\,z\in\mathbf{Z}_{j}.
\end{eqnarray}
Now, we apply Lemma \ref{LL33} to deduce that
\begin{equation}
\overline{q^{1}_{j,n}}-\overline{{q}^{2}_{j,n}}+(-1)^{n}\frac{l_{j}}{n\pi}(q^{1}_{j,n}\overline{q^{1}_{j,n}}-q^{2}_{j,n}\overline{q^{2}_{j,n}})
+q^{1}_{j,n}-q^{2}_{j,n}+(-1)^{n}(q^{1}_{j,n}-q^{2}_{j,n})=0.\label{3.16}
\end{equation}
For odd $n$,~(\ref{3.16}) is deduced to be 
\begin{equation}
\overline{q^{1}_{j,n}}-\overline{{q}^{2}_{j,n}}+(-1)^{n}\frac{l_{j}}{n\pi}(q^{1}_{j,n}\overline{q^{1}_{j,n}}-q^{2}_{j,n}\overline{q^{2}_{j,n}})=0.
\end{equation}
Therefore, we obtain $\Im \{\overline{q^{1}_{j,n}}-\overline{{q}^{2}_{j,n}}\}=0$.
Using $$q_{j,n}=\frac{2}{l_{j}}\int_{0}^{l_{j}}\Re\{q_{j}(x)\}\sin[\frac{n\pi}{l_{j}} (l_{j}-x)]dx+i\frac{2}{l_{j}}\int_{0}^{l_{j}}\Im\{q_{j}(x)\}\sin[\frac{n\pi}{l_{j}} (l_{j}-x)]dx.$$
Hence, we deduce $\Im\{q_{j}^{1}(x)-q_{j}^{2}(x)\}=0$. For even $n$, we obtain that $3\Re\{\overline{q^{1}_{j,n}}-\overline{{q}^{2}_{j,n}}\}+i\Im\{\overline{q^{1}_{j,n}}-\overline{{q}^{2}_{j,n}}\}$ is real, so  $\Im\{q_{j}^{1}(x)-q_{j}^{2}(x)\}\equiv0$. 
Using~(\ref{2.1})
\begin{eqnarray}\label{3.18}
\mathcal{F}(q^{1}_{j}-q^{2}_{j})(z)=\frac{2}{l_{j}}\int_{0}^{l_{j}}(q^{1}_{j}(x)-q^{2}_{j}(x))\sin[\frac{z\pi}{l_{j}} (l_{j}-x)]dx.
\end{eqnarray}
Hence, $\mathcal{F}(q^{1}_{j}-q^{2}_{j})$ is real on the real line. Using Cauchy-Riemann equation and Mean Value Theorem, $\mathcal{F}(q^{1}_{j}-q^{2}_{j})(z)$ is a constant on the real line. The constant must zero considering~(\ref{3.18}) at $z=0$. Therefore, the theorem is proven.



%Now, we go back to~(\ref{3.8}) and, then~(\ref{3.16}) to deduce
%\begin{equation}
%(q^{1}_{j,n}\overline{q^{1}_{j,n}}-q^{2}_{j,n}\overline{q^{2}_{j,n}})
%+\frac{n\pi}{l_{j}}[2(-1)^{n}+1](q^{1}_{j,n}-q^{2}_{j,n})=0,\label{3.18}
%\end{equation}
%That is, 
%\begin{equation}
%(q^{1}_{j,n}\overline{q^{1}_{j,n}}-q^{2}_{j,n}\overline{q^{2}_{j,n}})
%=-\frac{n\pi}{l_{j}}[2(-1)^{n}+1](q^{1}_{j,n}-q^{2}_{j,n}).\label{3.19}
%\end{equation}
%Since $\{q^{1}_{j,n}-q^{2}_{j,n}\}_{n}\in l^{2}$, we have $o((q^{1}_{j,n}-q^{2}_{j,n})^{2})=o(\frac{1}{n})$.
%Therefore, we deduce from~(\ref{3.19}) that
%\begin{equation}
%o(q^{1}_{j,n}\overline{q^{1}_{j,n}}-q^{2}_{j,n}\overline{q^{2}_{j,n}})
%=no(q^{1}_{j,n}-q^{2}_{j,n})=\sqrt{n}o(1).
%\end{equation}


\end{proof}
%\section{Appendix}
%For reader's convenience, we include the following classical lemma proved by E. C. Titchmarsh, and a modern proof can be found in \cite{Tang}.
%\begin{lemma}[Titchmarsh]\label{L2.1}
%Let $u\in\mathcal{E}'(\mathbb{R})$, then
%$$N_{\mathcal{F}(u)}(r)= \frac{|\mbox{ch supp } u|}{\pi}\big(r + o(1)\big)$$ and $$N_{f}(r)=\sum_{|z|\leq r}\frac{1}{2\pi}\oint_{z}\frac{f'(\omega)}{f(\omega)}d\omega,$$
%and the phrase $''\mbox{ch supp }u''$ means the covex hull of the effective support of $u$. Moreover, $N_{f}(r)$ is the counting function of the zeros of $f$ inside a ball of radius $r$. Let us count the zeros according to their multiplicities.
%\end{lemma}
%\begin{definition}
%We denote the following quantity as the zero density of an entire function $f$ of finite type.
%\begin{equation}
%\delta(f):=\lim_{r\rightarrow\infty}\frac{N_{f}(r)}{r}.
%\end{equation}
%For set $S$, we define similarly that
%\begin{equation}
%\delta(S):=\lim_{r\rightarrow\infty}\frac{N_{S}(r)}{r},
%\end{equation}
%where $N_{S}(r)$ is the number of the zeros of $S$ inside a ball of radius $r$ in $\mathbb{C}$.
%\end{definition}
%\begin{lemma}\label{L4.3}
%Let $f$, $g$ be two entire functions of finite type with density functions $\delta(f)$ and
%$\delta(g)$ respectively. Then, 
%\begin{equation}\nonumber
%\delta(fg)=\delta(f)+\delta(g),
%\end{equation}
%and
%\begin{equation}
%\delta(f+g)=\max\{\delta(f),\delta(g)\},\label{220}
%\end{equation}if the types of two functions are not equal.
%\end{lemma}
%\begin{proof}
%We refer to B. Ya. Levin's books \cite[p.\,52]{Levin,Levin2} for a detailed introduction.  It is all about completely regular growth functions.


%\end{proof}
%\begin{lemma}\label{L4.4}
%For $f(x)$ in $L^{2}(0,\pi)$, $z\in\mathbb{C}$, and $\int_{0}^{\pi}e^{-iz x}f(x)dx=\int_{0}^{\pi}\cos{z x}f(x)dx-i\int_{0}^{\pi}\sin{z x}f(x)dx$, the entire function $$\int_{0}^{\pi}e^{-iz x}f(x)dx;\\ \int_{0}^{\pi}\cos{z x}f(x)dx,\mbox{ and }\int_{0}^{\pi}\sin{z x}f(x)dx$$ have the same zero density in $\mathbb{C}$.
%\end{lemma}
%\begin{proof}
%Let us simply apply Titchmarsh's results \cite{Titchmarsh}.

%\end{proof}

\begin{acknowledgement}
The author ingenuously appreciates the research funding supported by NSTC under the project number 113-2115-M-131 -001. The content of this manuscript does not necessarily reflect the position or the policy of the administration, and no official endorsement should be inferred, neither.
\end{acknowledgement}

\begin{thebibliography}{widest-label}
\bibitem{Alb}Albeverio S., Hryniv R., and Nizhnik L., Inverse spectral problems for nonlocal Sturm-Liouville operators, Inverse Problems, 23, 523--535 (2007).
\bibitem{Alb2}Albeverio S. and Nizhnik L. P., Schr\"{o}dinger operators with nonlocal point interactions, J. Math. Anal. Appl., 332, 884--95 (2007 ).


\bibitem{Belishev}Belishev M. I., Boundary spectral inverse problem on a class of graphs (trees) by the BC method, Inverse Problems, 20, 647--672 (2004).
%\bibitem{Boman}Boman J. and Kurasov P., Symmetries of quantum graphs and the inverse scattering problem, Adv. Appl. Math., 35, 58--70 (2005).
%\bibitem{Bon2}Bondarenko N.P., Buterin S. A., Vasiliev S. V., An inverse spectral problem for Sturm--Liouville operators with frozen
%argument, J. Math. Anal. Appl., 472, 1028--1041 (2019).
\bibitem{Bon}Bondarenko N. P., A partial inverse problem for the Sturm--Liouville operator on a star-shaped graph, Anal. Math.
Phys., no. 8, 155--168 (2018).
%\bibitem{Bru}Br\"{u}ning J. and Geyler  V. A., Scattering on compact manifolds with infinitely thin horns, J. Math. Phys., 44 , 371--405 (2003).
\bibitem{Carlson}Carlson R., Inverse eigenvalue problems on directed graphs, Trans. Amer. Math. Soc., 351, 4069--4088 (1999).
\bibitem{Exner}Exner P.  and Post O., Convergence of spectra of graph-like thin manifold, J. Geom. Phys., 54, 77--115 (2005).
\bibitem{Freiling}Freiling G. and Yurko V. A., Inverse problems for differential operators on graphs with general matching conditions, Applicable Analysis, 86, no. 6, 653--667 (2007).
\bibitem{Gerasimenko}Gerasimenko N. I. and Pavlov B. S., Scattering problems on non-compact graphs, Theor. Math. Phys., 74, 230--240  (1988).

\bibitem{Kostrykin}Kostrykin V. and Schrader R., Kirchhoff's rule for quantum wires. II. The inverse problem with possible applications to quantum computers, Fortschr. Phys., 48, 703--716 (2000).
\bibitem{Kottos}Kottos T. and Smilansky U., Quantum graphs: A simple model for chaotic scattering, J. Phys. A: Math. Gen., 36, 3501--3524  (2003).
%\bibitem{Kravchenko}Kravchenko K. V., Diff. Eqns 36 517--23, 2000.  (Engl. transl.)

\bibitem{Kuchment}Kuchment P., Quantum graphs. I. Some basic structures, Waves in Random Media, 14, 107--128 (2004).
\bibitem{Kurasov}Kurasov P. and Novaszek M., Inverse spectral problem for quantum graphs, J. Phys. A: Math. Gen., 38, 4901--4915 (2005).
\bibitem{Kurasov2}Kurasov, P., Spectral Geometry of Graphs, Operator Theory: Advances and Applications
Volume 293, Birkh\"{a}user, Berlin,1979.
%\bibitem{Levin}Levin B. Ja., Distribution of Zeros of Entire
%Funtions, revised edition, Translations of Mathematical
%Monographs, American Mathematical Society, Providence, 1972.
\bibitem{Levin2}Levin B. Ja., Lectures on Entire Functions,
Translation of Mathematical Monographs, V. 150, AMS, Providence, 1996.

%\bibitem{Marchenko} Marchenko V. A., Sturm-Liouville Operators and Their Applications, Naukova Dumka, Kiev, 1977. (Russian); English transl. Birkh\''{a}user Verlag, Basel, 1986.
\bibitem{Naimark}Naimark K. and Solomyak M., Eigenvalue estimates for the weighted Laplacian on metric trees, Proc. London Math. Soc., 80, 690--724  (2000).


\bibitem{Nizhnik1}Nizhnik L. P., Inverse eigenvalue problems for nonlocal Sturm--Liouville operators on a star graph, Methods of Functional Analysis and Topology,
Vol.18, no. 1, 68--78. (2012)
\bibitem{Nizhnik2}Nizhnik L. P., Inverse nonlocal Sturm-Liouville problem, Inverse Problems,  26, 125006 (2010). 


\bibitem{Novikov}Novikov S. P., Schr\"{o}dinger operators on graphs and symplectic geometry, E. Bierstone, B. Khesin, A. Khovanskii, and J. Marsden (Eds.), The Arnoldfest, Proceedings of a Conference in Honour V. I. Arnold for His Sixtieth Birthday, Fields Institute Communications Series, Vol. 24, Amer. Math. Soc., Providence, RI, 397--413 (1999).
\bibitem{Pivovarchik1}Pivovarchik V., Inverse problem for the Sturm-Liouville equation on a simple graph, SIAM J. Math. Anal., 32, 801--819 (2001).
%\bibitem{Pivovarchik2}Pivovarchik V., Inverse problem for the Sturm-Liouville equation on a star-shaped graph, Math. Nachr.,  280, no. 13. 14., 1595--1619  (2007).
\bibitem{Pokornyi}Pokornyi Yu. and Pryadiev  V., The qualitative Sturm-Liouville theory on spatial networks, J. Math. Sci. (N.Y.) 119, no. 6, 788--835 (2004).
%\bibitem{Tang}Tang Siu-Hung and Zworski M., Potential scattering on the real line, Department of mathematics, U.C. Berkeley, http://math.berkeley.edu/~zworski/tz1.pdf.
%\bibitem{Titchmarsh}Titchmarsh E. C., The zeros of certain integral functions, Proceedings of The London Mathematical Society, 283--302 (1926).
\bibitem{Yang}Yang, C. -F., Inverse spectral problems for the Sturm--Liouville operator on a d-star graph, J. Math. Anal. Appl., 365, 742--749  (2010).
\bibitem{Yurko1}Yurko V. A., Inverse spectral problems for Sturm-Liouville operators on graphs, Inverse Problems, 21, 1075--1086 (2005).
\bibitem{Yurko2}Yurko V. A., An inverse problem for higher-order differential operators on star-type graphs, Inverse Problems, 23, 893--903 (2007).
\bibitem{Shieh}Wang, Y. and Shieh, C. -T., Inverse problems for Sturm--Liouville operators on a star-shaped graph with mixed spectral data, Applicable Analysis, 1--10 (2019).




\end{thebibliography}
\end{document}
